\chapter*{}

\vspace{-4cm}

\textsf{Este libro está enfocado a aquellos alumnos que van a emprender carreras científicas que necesitan de más contenidos matemáticos que los aprendidos en bachillerato, por tanto, se suponen conocidos todos los temas allí tratados, quien lea el libro debe saber manejarse con soltura ante sistemas de ecuaciones lineales, matrices y determinantes, conocer el producto escalar y vectorial de vectores y saber derivar e integrar con soltura.} 

\textsf{Se mostrarán, de forma sencilla y con ejercicios resueltos (y propuestos con solución) algunos de los temas que se afrontarán en los primeros cursos de carrera. Para ello es necesario perder el rigor necesario en matemáticas pero será en pos de aclarar los conceptos, al cursar estudios universitarios se recuperará el rigor que ahora, de momento, abandonamos. Espero que sea de utilidad.}

\textsf{He incluidos los temas que figuran en el índice y considero este trabaja abierto a posibles modificaciones e inclusión de otros temas de interés.}


\vspace{1cm}

\textbf{Estructura}

No he sido demasiado exigente a la hora de usar el mismo formato en todos los textos, pero aparecerán:

\begin{cuadro-gris}
	
\end{cuadro-gris}

\begin{cuadro-naranja}
	
\end{cuadro-naranja}

\begin{destacado}
	
\end{destacado}

\begin{miejercicio}
	
\end{miejercicio}

\begin{mipropuesto}
	
\end{mipropuesto}

\begin{theorem}
	
\end{theorem}

\begin{definition}
	
\end{definition}







\begin{comment}

%%%%%%%%%%%%%%%%%%%%%%%%%%%%%%%%%%%. SECCIONES
\chapter{texto}
\begin{tikzpicture}
	\fill [left color=red!50, right color=teal!50] (0,0) rectangle (6.5,.2);
	\fill [left color=teal!50, right color=blue!50] (6.5,0) rectangle (11.5,.2);
	\end{tikzpicture}

\vspace{1cm}
\section{texto}
\begin{tikzpicture}
	\fill [left color=red!50, right color=teal!50] (0,0) rectangle (3.5,.1);
	\fill [left color=teal!50, right color=blue!50] (3.5,0) rectangle (7.5,.1);
	\end{tikzpicture}
\vspace{0.5cm}

\subsection{texto}
\begin{tikzpicture}
	\fill [left color=red!50, right color=teal!50] (0,0) rectangle (3.5,.01);
	\fill [left color=teal!50, right color=blue!50] (3.5,0) rectangle (7.5,.01);
	\end{tikzpicture}
\vspace{0.5cm}


%%%%%%%%%%%%%%%%%%%%%%%%%%%%%%%%%%%. \begin{ ------>. 
detsacado;  cuadro-naranja;  cuadro-gris;  miejercicio (solución extensa);  mipropuesto (solución corta y fuera del cuadro)

%%%%%%%%%%%%%%%%%%%%%%%%%%%%%%%%%%%. CURIOSIDAD
\vspace{1cm}
\color{ForestGreen!80}
\rule{250pt}{0.2pt}
Texto
\vspace{-8mm}
\begin{flushright}
\rule{250pt}{0.2pt}		
\end{flushright}	
\color{black}
\end{comment}



\emph{\normalsize{Este} documento se comparte bajo licencia `Creative Commons Attribution-NonCommercial-ShareAlike (4.0 CC BY-NC-SA)'}


\begin{multicols}{2}
\begin{figure}[H]
	\centering
	\includegraphics[width=.5
	\textwidth]{imagenes/licencia.png}
	
\end{figure}
\begin{figure}[H]
	\centering
	\includegraphics[width=.25
	\textwidth]{imagenes/firma.png}
\end{figure}
\end{multicols}

\newpage




\newpage

$\,$

\newpage

$\,$

\vspace{1cm} 

$\,$

\begin{figure}[H]
	\centering
	\includegraphics[width=.80\textwidth]{imagenes/presenta1.png}
\end{figure}


\begin{center}

\Huge{\textbf{Ampliación de Matemáticas:}}

\huge{\textbf{del Instituto a la Facultad}} 


\Large{\textbf{\textit{Una visión introductoria}}}

\vspace{10mm}
\begin{figure}[H]
	\centering
	\includegraphics[width=.78\textwidth]{imagenes/presenta2.png}
\end{figure}

\vspace{2cm}
\begin{flushright}
	\normalsize{\emph{Ignacio Vallés Oriola}}
\end{flushright}


\end{center}





